\section{Discussion}
The two finite-collar fermionic quantities weight the angular spectrum in different ways. Changing the Lorentzian preparation across one crossing measures an integer multiplicity in a single shell. Changing the charge in the compact modulus measures a continued signed sum over \emph{all} shells. Table~\ref{tab:eta} gives examples realizing either phenomenon without the other. A mirror or change of spin lift may exchange the conclusion even at the same radius and fixed $|P|$. A vector-like field has a vanishing vector-number difference although its insertion rule and its modulus can still be nontrivial.

Three limiting operations enter these results. Bringing $h$ to zero with a continuous elliptic domain chooses a definite rule for a contacting shell. Moving a preparation surface to the throat at the same time can change the number measured by that state experiment. Taking $R$ to infinity after the compact relative determinant yields the same formula as a shell-first Jost threshold with angular Abel summation, for the fifteen lifts under study. This agreement does not establish that the two limiting operations commute in a full scattering trace. Individual absolute one-loop moduli would require control of branch-even high-frequency terms.

The finite-boundary bosonic negative directions also depend on which boundary data fluctuate. E1 closes the harmonic flux direction. E2 integrates one real collective coordinate and, for a sufficiently short collar, releases an endpoint coordinate in some scalar harmonics. In the strict infinite-boundary action the collective flux coefficient vanishes. This is consistent with a negative coefficient at every finite radius. A nonnegative bulk integrand suffices for E1; it settles neither the sign of an E2 boundary term nor the contour of a complex scalar path integral. The connected-sum example supplies an independent separation between an asymmetric eta invariant and actual throat contact.

The common underlying calculation uses a test field on a prescribed solution and specified boundary data. The bosonic action, Morse index, Fock insertion number, modulus of the relative Dirac determinant and local phase have distinct definitions. Turning their values into branch probabilities would require a gravitational integration cycle, a measure, a prescription for the state of the other fields and, in sectors with fermionic zero modes, specified insertions. Global noncompact determinants, nonlinear stability and backreaction are outside the present boundary problem.
